\section{Uniform counterfactual inference and multiplier bands}\label{sec:uniform}

Let $T_\tau(\vartheta)$ denote a CDF, quantile, tail, or effect target. The map
\[
 \vartheta\mapsto T(\vartheta)\in C(\cT)
\]
is uniformly Fr\'echet differentiable on the compact target class. Because the estimated factual model has two parameters, the limiting processes are finite rank.

For a fixed absolute-intervention CDF,
\[
 F_{a,w,z}^{\cf}(\sigma)
 =\Phi\left(\frac{z-m_a(w)}{\sigma\sqrt{V_a(w)}}\right),
\]
and
\[
 d_\sigma^{\cf}(a,w,z)
 =-\frac{c_{a,w,z}\phi(c_{a,w,z})}{\sigma_0}.
\]

\begin{theorem}[Uniform mixed-rate target processes]\label{thm:processes}
Under \Cref{ass:core,ass:regime} and compactness of $\cT$,
\[
 \sqrt{Nn}\left(\widehat F^{\cf}-F^{\cf}\right)
 \Rightarrow
 \mathbb G_{\cf}
 \quad\text{in }C(\cT),
\]
where
\[
 \mathbb G_{\cf}(a,w,z)=d_\sigma^{\cf}(a,w,z)Z_\sigma,
 \qquad
 Z_\sigma\sim N(0,\sigma_0^2/2).
\]
Factual CDFs satisfy
\[
 \sqrt{\Lambda_N}(\widehat F^{\fact}-F^{\fact})
 \Rightarrow d_\theta^{\fact}(w,z)Z_\theta,
 \qquad
 Z_\theta\sim N(0,2\theta_0/T).
\]
For the absolute-intervention distributional effect
\[
 \Delta_{a,w,z}=F_{a,w,z}^{\cf}-F_{w,z}^{\fact},
\]
the slow first-order limit is
\[
 \sqrt{\Lambda_N}(\widehat\Delta-\Delta)
 \Rightarrow-d_\theta^{\fact}(w,z)Z_\theta.
\]
\end{theorem}

\begin{theorem}[Uniform quantile and tail processes]\label{thm:quantile-tail}
Let $\mathcal P\Subset(0,1)$ be compact and let
$Q_{a,w}(p)=m_a(w)+\sigma\sqrt{V_a(w)}\,\Phi^{-1}(p)$ be the
fixed-intervention quantile surface. Under the conditions of
\Cref{thm:processes},
\[
 \sqrt{Nn}\big(\widehat Q-Q\big)
 \Rightarrow\sqrt{V_a(w)}\,\Phi^{-1}(p)\,Z_\sigma
 \qquad\text{in }C(\widetilde\Theta\times\cW\times\mathcal P),
\]
and the fixed-intervention tail process
$\sqrt{Nn}(\widehat S-S)$, $S:=1-F^{\cf}$, converges to
$-\mathbb G_{\cf}$. Factual and distributional-effect versions hold at rate
$\sqrt{\Lambda_N}$ with the corresponding $\theta$-loadings. Proofs are in
\ifdefined\MERGEDDOC Part~S4 (\Cref{mpart:wp4})\else part S4 of the technical supplement\fi.
\end{theorem}

\begin{remark}[Counterfactual-distribution parallel]\label{rem:cfm}
The objects of \Cref{thm:processes,thm:quantile-tail} are the dynamic-field analogues of counterfactual distribution and quantile processes in microeconometrics \citep{ChernozhukovFernandezValMelly2013}. The structural difference is the limit geometry: a two-parameter factual experiment makes every target process finite rank, so simultaneous inference reduces to a two-dimensional Gaussian draw (\Cref{thm:bands}) rather than to empirical-process approximations over Donsker classes.
\end{remark}

The slow effect limit is constant in the absolute intervention coordinate because the counterfactual part is learned at the faster scale rate. A relative intervention $a=\eta\theta$ yields intervention-indexed slow uncertainty when that is scientifically appropriate.

For feasible inference, let $S_{k,N,n}(\widehat\vartheta)$ be the complete modal score block and set
\[
 \widehat B_k=A_{N,n}^{-1}S_{k,N,n}(\widehat\vartheta),
 \qquad
 \widehat J_B=\sum_{k=1}^N\widehat B_k\widehat B_k^\top.
\]
With independent Gaussian multipliers $\xi_k$ define
\[
 H^*=\widehat J_B^{-1}\sum_{k=1}^N\xi_k\widehat B_k.
\]
Conditionally on the data,
\[
 H^*\sim N(0,\widehat J_B^{-1}).
\]
Thus the score multiplier is equivalent to a direct two-dimensional Gaussian draw with estimated inverse-information covariance.

\begin{theorem}[Simultaneous confidence bands]\label{thm:bands}
The multiplier target processes obtained by multiplying the estimated derivative surfaces by the corresponding coordinates of $H^*$ converge conditionally to the processes in \Cref{thm:processes}. Conditional sup-norm quantiles yield asymptotically valid simultaneous bands.

At derivative zeros, define
\[
 \widehat s_\kappa(\tau)
 =\max\{\widehat s(\tau),\kappa_{N,n}\norm{\widehat s}_\infty\}.
\]
For fast targets require
\[
 \kappa_{N,n}\downarrow0,
 \qquad
 \sqrt{Nn}\,\kappa_{N,n}\to\infty.
\]
For slow targets also require
\[
 \sqrt{\Lambda_N}\,\kappa_{N,n}\to\infty,
 \qquad
 \frac{\sqrt{\Lambda_N}}{\sqrt{Nn}\,\kappa_{N,n}}\to0.
\]
The resulting floor-studentised bands remain valid uniformly, including at first-order derivative zeros.
\end{theorem}

\begin{figure}[t]
\centering
\resizebox{.8\textwidth}{!}{% Schematic: floor-studentised simultaneous band and estimated contact set.
% Effect surface tau -> Delta(tau) (one-sided null Delta <= 0), band of
% half-width c* s_kappa / r, floor active near the derivative zero, and the
% estimated contact set where the studentised surface exceeds -eta_{N,n}.
\begin{tikzpicture}[font=\small, scale=1.0,
  ax/.style={-{Stealth[length=2.2mm]}, thick},
  lbl/.style={font=\footnotesize\itshape, text=black!70}]
% effect curve Delta(tau) = -(tau-3)^2/6 + 0 : contact at tau=3
\pgfmathsetmacro{\shift}{2.2}
\draw[ax] (0,\shift) -- (8.6,\shift) node[below left=1pt]{$\tau$};
\draw[ax] (0,0) -- (0,4.4);
\node[anchor=south, rotate=90] at (-.6,2.2) {$\widehat\Delta(\tau)$, band};

% band (drawn first, behind): half-width h(tau) = .30 + .22*exp(-((tau-3)^2))  floor bump near zero-derivative point tau=3?
% use s(tau) ~ |tau-3|/3 studentiser with floor kappa: h = max(.14, .38*|tau-3|/3) + .12
\fill[blue!10]
  plot[domain=0.3:8.2, samples=80]
    (\x,{\shift-((\x-3)^2)/6 + max(0.16,0.36*abs(\x-3)/3)+0.10})
  -- plot[domain=8.2:0.3, samples=80]
    (\x,{\shift-((\x-3)^2)/6 - max(0.16,0.36*abs(\x-3)/3)-0.10});
% true effect
\draw[thick, black!70, dashed] plot[domain=0.3:8.2, samples=80]
  (\x,{\shift-((\x-3)^2)/6});
% estimate (slightly wiggly)
\draw[thick, blue!60!black] plot[domain=0.3:8.2, samples=90]
  (\x,{\shift-((\x-3)^2)/6+0.09*sin(160*\x)});
% zero line
\draw[black!50] (0,\shift) -- (8.6,\shift);
\node[font=\scriptsize, black!50, anchor=west] at (8.15,\shift+0.16) {$0$};

% floor-active region: where .36|tau-3|/3 < .16  => |tau-3| < 1.333
\draw[decorate, decoration={brace, mirror, amplitude=4pt}, black!60]
  (1.667,0.28) -- (4.333,0.28)
  node[midway, below=5pt, lbl]{floor active: $\widehat s<\kappa_{N,n}\norm{\widehat s}_\infty$};
% contact set (true): tau = 3
\fill[red!70!black] (3,\shift) circle (2pt);
\node[lbl, red!70!black, above right=1pt and 2pt] at (3,\shift) {$\cC_0$};
% estimated contact set: where r Delta / s_kappa >= -eta: interval around 3
\draw[very thick, red!60!black] (2.05,0.72) -- (3.95,0.72)
  node[midway, below=2pt, lbl, red!60!black]
  {$\widehat{\cC}_{N,n}$ (shrinks to $\cC_0$ at rate driven by $\eta_{N,n}/r$)};
% band label
\node[lbl, blue!50!black, align=left] at (6.7,3.75)
  {simultaneous band\\ $\widehat T\pm c^*_{1-\alpha}\widehat s_\kappa/r$};
\draw[-{Stealth[length=2mm]}, black!50] (6.05,3.55) -- (5.35,3.05);
\end{tikzpicture}
}
\caption{Schematic of the floor-studentised simultaneous band and the estimated contact set for a one-sided effect hypothesis. Near the derivative zero the floor $\kappa_{N,n}\norm{\widehat s}_\infty$ keeps the band width above the quadratic remainder; the estimated contact set $\widehat{\cC}_{N,n}$ retains the true contact set and shrinks onto it because $\eta_{N,n}\to\infty$ while $\eta_{N,n}/r_{N,n}\to0$.}
\label{fig:band-contact}
\end{figure}

Equality hypotheses use sup-norm or integrated-square statistics. For one-sided hypotheses $G(\tau)\le0$, define the estimated contact set
\[
 \widehat{\mathcal C}_{N,n}
 =\left\{\tau:
 \frac{r_{N,n}\widehat G(\tau)}{\widehat s_\kappa(\tau)}
 \ge-\eta_{N,n}
 \right\},
\]
with $\eta_{N,n}\to\infty$ and $\eta_{N,n}/r_{N,n}\to0$. Bootstrapping the positive-part supremum over this set gives asymptotically valid and consistent stochastic-dominance, tail-sign, and quantile-sign tests, following the logic of moment-selection procedures \citep{AndrewsSoares2010,LintonSongWhang2010,FangSantos2019}.
